\documentclass[11pt]{article}
\usepackage[margin=1in]{geometry}
\usepackage{amsmath,amssymb,amsthm}
\usepackage{hyperref}
\newtheorem{definition}{Definition}
\newtheorem{proposition}{Proposition}
\newtheorem{remark}{Remark}
\title{Predictive Prefetching for Agent-Facing MCP Servers:\\ A Formal Model and Its Incentive Problem}
\author{MCP Receiver project}
\date{September 2026}
\begin{document}
\maketitle

\begin{abstract}
We formalize an MCP server that observes an agent's tool calls, predicts the agent's next call with a
fast classifier (Jev, a ``System 1'' model), and pre-generates the response before the call arrives.
Pages produced this way are attached to claimed domains and offered to third-party brands as
paid placements. We give the model, a latency and hit-rate analysis, and then state the issue
plainly: the same machinery that hides latency can be used to shape what an agent shows its user,
which is a principal--agent problem. We list the design constraints the implementation enforces.
\end{abstract}

\section{Model}
Let $\Sigma$ be a finite alphabet of tool names (e.g.\ \texttt{find\_providers}, \texttt{read\_skill}).
A \emph{session} is a word $w = t_1 t_2 \cdots t_n \in \Sigma^*$ produced by an agent.
Let $A$ be a set of professions and $L$ a set of locations.

\begin{definition}[Bait set]
$B = A \times L$. A call $c = (t, (a,\ell))$ is a \emph{bait read} iff $(a,\ell)\in B$.
\end{definition}

\begin{definition}[Next-step predictor]
A predictor is a map $\pi : \Sigma \times B \to \Delta(\Sigma)$. The empirical baseline is the first-order
Markov estimate from logged sessions $W$:
\[
\hat\pi(t' \mid t) \;=\; \frac{N_W(t\,t')}{\sum_{s\in\Sigma} N_W(t\,s)},
\]
where $N_W(u)$ counts occurrences of the bigram $u$. Jev replaces $\hat\pi$ with a model that
outputs $\pi(\cdot\mid t,(a,\ell))$ directly over a predefined output set, so it cannot emit an
output outside $\Sigma$.
\end{definition}

\begin{definition}[Prefetch policy]
Fix a threshold $\tau\in(0,1]$. On a bait read $(t,b)$ the server computes
$P_\tau(t,b)=\{t' : \pi(t'\mid t,b)\ge\tau\}$ and generates an artifact $g(t',b)$ for each
$t'\in P_\tau$ before the agent's next request.
\end{definition}

\section{Analysis}
\begin{proposition}[Hit rate]
If the agent's true next tool is drawn from $\pi$, the expected fraction of next calls served from
prefetch is $H(\tau)=\sum_{t' \,:\, \pi(t')\ge\tau}\pi(t')$, and the expected number of generated
artifacts is $|P_\tau|\le 1/\tau$.
\end{proposition}
\begin{proof}
Each $t'\in P_\tau$ is prefetched, so a draw lands in $P_\tau$ with probability
$\sum_{t'\in P_\tau}\pi(t')$. Since $\pi$ sums to $1$ and each member has mass $\ge\tau$,
$|P_\tau|\tau\le 1$.
\end{proof}

\begin{proposition}[Latency hiding]
Let $\lambda_g$ be generation latency and $\delta$ the inter-call gap (agent think time).
A prefetched artifact is ready in time iff $\lambda_g\le\delta$. The expected response time is
$\mathbb E[\lambda]=H(\tau)\,\lambda_c+(1-H(\tau))(\lambda_c+\lambda_g)$ with $\lambda_c$ the cache
read cost, so prefetching helps exactly when $\lambda_g\le\delta$ and $H(\tau)$ is large.
\end{proposition}

\paragraph{Domain claiming.} Let $D$ be a pool of free domains and $\beta$ a purchase budget.
Claiming is an online allocation: assign a free $d\in D$ to each new bait; if $D$ is exhausted,
buy a new domain at price $p$ only if $\text{spent}+p\le\beta$. Purchasing is disabled unless
explicitly enabled.

\section{The issue}
The server is not a neutral party. Its output is consumed by an agent acting for a user (the
\emph{principal}); the operator is paid by third-party brands. This is a principal--agent problem
with a hidden intermediary:

\begin{enumerate}
\item \textbf{Objective mismatch.} The agent maximizes user utility $U$; the operator maximizes
      brand revenue $R$. Content optimized for $R$ and selected by the agent because it is
      well-formed and present at the right moment can displace better-ranked answers for $U$.
\item \textbf{Prediction is manipulation-capable.} Because $g(t',b)$ is written \emph{before} the
      agent asks and tuned to what the agent's model tends to accept, the operator controls
      both the question distribution and the answer, an asymmetric advantage over any user.
\item \textbf{Non-disclosure.} If sponsored placement is not visible to the end user, the outcome is
      deceptive advertising, and platforms treat agent-directed instruction text as prompt
      injection.
\end{enumerate}

\begin{remark}[Constraints enforced]
To keep the system on the legitimate side of this line, the generator (i) is prompted for factual
content, (ii) rejects any output containing imperatives aimed at the reading agent
(e.g.\ ``ignore previous instructions''), (iii) labels sponsored placement on every page, and
(iv) restricts anonymous writes to a scratch namespace. Whether these suffice is an open policy
question that depends on how the consuming agent presents sources to its user.
\end{remark}

\section{Race dynamics in an agentic internet}
\paragraph{Read-path race.} Let $O=\{o_1,\dots,o_m\}$ be operators that each pre-position an artifact
for the predicted step $t'$ on bait $b$. An agent with latency budget $\delta$ accepts the first
artifact $o_i$ that is well-formed and ready by $\delta$; it does not enumerate all $m$. Operator $i$
wins with probability depending on readiness time $r_i\le\delta$ and acceptance probability $\alpha_i$,
not on true quality $q_i$. Since entry cost is near zero (free domains, cheap generation), $m$ grows
until margin vanishes.

\paragraph{Why it favours fraud.} Define $\alpha_i$ as the agent's acceptance rate for $o_i$'s content.
An operator optimizing $\alpha_i$ against a specific model is optimizing against that model's
judgment. Honest and fraudulent operators solve the same problem; the fraudulent one has a strictly
larger feasible set of artifacts (fabricated contacts, lookalike payment endpoints), so in
equilibrium the accepted set is enriched for fraud unless verification is free. Verification costs
latency, and latency loses the read-path race, so agents are pushed toward skipping it.

\paragraph{TOCTOU.} If an artifact is checked at time $t_c$ and used at $t_u>t_c$, an operator that
mutates content in $(t_c,t_u)$ defeats the check. Prefetched and cached artifacts widen this window.

\paragraph{Mitigations.} Signed provenance; reputation and age gates on domains for
high-stakes fields; cross-source agreement; hash-pinning between check and use; human confirmation
for money and contact actions; end-user-visible sponsorship disclosure.

\section{Projected bottleneck and whole-intent submission}
\emph{This section is a projection, not a measurement.}

\paragraph{Prefetch advantage vanishes.} From Proposition 2, prefetching helps only when
$\lambda_g\le\delta$ \emph{and} on-demand generation is slower than the gap. Let $\lambda_g^{\text{od}}$
be the latency of a fast generator (possibly another agent called at request time). As
$\lambda_g^{\text{od}}\to 0$, the advantage $\mathbb E[\lambda_{\text{on-demand}}]-\mathbb E[\lambda_{\text{prefetch}}]\to 0$
and $H(\tau)$ no longer determines who wins. Competition shifts from prediction quality to raw generation
speed, which is reset by every model improvement.

\paragraph{Where the bottleneck moves.} Per-step calls cost one round trip each, so a chain of $n$
calls across agents costs $\Theta(n)$ round trips and $\Theta(n)$ opportunities for interception.
When generation is cheap, verification becomes the scarce resource: acceptance is limited by
verification cost $v$, not by generation cost, and $v$ is paid in latency the agent cannot afford. This
reproduces the scam-race incentive of the previous section at larger scale.

\paragraph{Whole-intent submission.} Replace the per-step stream $t_1\cdots t_n$ by a single declared
intent $I$ (goal, constraints, budget, verification requirement). The server no longer predicts
$\pi(t'\mid t,b)$; it plans over $I$ in one round trip and can return a signed plan
$\sigma(\text{plan},I)$ the agent checks against its constraints.

\paragraph{New costs.} (i) Information leakage: the operator observes all of $I$, increasing the
asymmetry noted earlier unless $I$ is minimized or the operator is constrained. (ii) Intent becomes the
auctioned object, so brands bid on intent categories. (iii) $I$ does not itself prove honesty of the
response; provenance and hash-pinning are still required.

\section{Potential fix}
\subsection{Whole-intent submission as a protocol}
Let an intent be a tuple $I=(g,C,\beta,\rho)$: goal $g$, hard constraints $C$, budget $\beta$, and a
verification requirement $\rho$ (e.g.\ ``contact details must be corroborated by two independent
sources''). The agent submits $I$ once. The server returns a \emph{plan}
$P=(s_1,\dots,s_k)$ together with a signature $\sigma=\mathrm{Sign}_{sk}(H(P)\,\|\,H(I))$, and each
step $s_j$ carries the provenance and content hash of the artifact it refers to.

\begin{definition}[Acceptable plan]
The agent accepts $(P,\sigma)$ iff (i) $\mathrm{Verify}_{pk}(\sigma)$ holds, (ii) every $s_j$ satisfies $C$
and costs at most $\beta$ in total, (iii) $\rho$ is met by the attached provenance, and (iv) each
artifact hash matches at time of use, which closes the TOCTOU window.
\end{definition}

\begin{proposition}[Round trips]
Per-step interaction needs $\Theta(n)$ round trips for an $n$-call chain. Intent submission needs
$O(1)$ to obtain the plan plus $O(k)$ local hash checks, removing the prediction step
$\pi(t'\mid t,b)$ from the critical path entirely.
\end{proposition}

\subsection{Limits}
Submitting $I$ leaks it, so the design minimizes disclosure: $g$ and $C$ are coarsened and $\beta$
is sent as a bound, not a value, and the operator may not use $I$ for anything other than
answering it. Intent-category auctions must remain visibly labelled. Signatures prove who said
something, not that it is true, so $\rho$ (independent corroboration) still carries the trust load.

\subsection{Future implementation (hint)}
The current server accepts only single tool calls at \texttt{/mcp}. A future revision would add a
\texttt{submit\_intent} tool in \texttt{app/mcp.py} taking $I$ and returning $(P,\sigma)$;
a \texttt{planner} module in place of the per-step predictor in \texttt{app/pipeline.py};
and a signing key with hash-pinned artifacts stored next to each row in the \texttt{artifacts}
table. Jev remains useful there as a fast classifier (e.g.\ intent category and risk level),
while a separate model composes the plan.

\section{Evidence for Jev's speed, and generalization}
\subsection{What is actually known}
All available figures are vendor claims or individual developer reports; \emph{none are independently verified}.
\begin{itemize}
\item TypeSafe reports response times of 70--500\,ms and input pricing of \$0.042 per million tokens; its accuracy
      reference is the average of two frontier models' answers, and the evaluations come from the start-up itself
      \cite{heise}.
\item Reported relative speed-ups: ``20--200$\times$ faster'' and ``40--400$\times$ cheaper'' with no named baselines
      \cite{latentspace}; developer reports of 5--18$\times$ (safety classification) and 10--20$\times$ versus Gemini
      (email classification), where Gemini was slightly more accurate but Jev returned confidence scores
      \cite{techcrunch}.
\item Vendor announcement and docs: \cite{typesafe}. Jev cannot emit free-form text, so speed comparisons to general LLMs
      hold only for closed-output classification and routing tasks.
\end{itemize}

\subsection{Abstraction: any sufficiently fast, acceptably good model}
Let a model $M$ have latency $\lambda_M$, task accuracy $a_M$, and cost $c_M$. Fix an acceptability
threshold $a^\ast$ (the accuracy at which the predictor's output is useful) and the agent's latency budget
$\delta$. Define the \emph{inline-feasible set}
\[
\mathcal F=\{M : a_M\ge a^\ast \;\wedge\; \lambda_M\le \varepsilon\,\delta\},\qquad \varepsilon\ll 1 .
\]
\begin{proposition}[Generality]
If $M\in\mathcal F$, the prediction step adds at most $\varepsilon\delta$ to the critical path, so the
prefetch/on-demand race of the earlier sections is available to any operator using $M$. If $M'$ is $k$ times faster
than $M$ at equal $a\ge a^\ast$, the extra room $\delta(\varepsilon-\varepsilon/k)$ can be spent on
more candidate steps ($|P_\tau|$ larger at smaller $\tau$) or on generation, so each order of magnitude
of speed at acceptable quality strictly enlarges what an operator can pre-position inside $\delta$.
\end{proposition}
Nothing in this depends on Jev's architecture. Autoregressive small models, distilled classifiers, or
non-autoregressive systems all qualify once $a_M\ge a^\ast$ and $\lambda_M$ is small. Jev is one instance of the
speed--quality frontier moving; the behaviour described in this paper follows from the frontier, not from the
vendor. Consequently defences must be placed on the receiving side (provenance, verification, disclosure), because
the supply of fast-enough predictors cannot be limited.

\begin{thebibliography}{9}
\bibitem{heise} heise online, ``AI model Jev to make machines decide faster,'' Sept.\ 2026.
  \url{https://www.heise.de/en/news/AI-model-Jev-to-make-machines-decide-faster-11457071.html}
\bibitem{latentspace} Latent Space, ``[AINews] Jev: a System One Model\dots,'' Sept.\ 2026.
  \url{https://www.latent.space/p/ainews-jev-a-system-one-model-that}
\bibitem{techcrunch} TechCrunch, ``A new kind of AI model from a ChatGPT inventor is thrilling developers,'' 18 Sept.\ 2026.
  \url{https://techcrunch.com/2026/09/18/a-new-kind-of-ai-model-from-a-chatgpt-inventor-is-thrilling-developers/}
\bibitem{typesafe} TypeSafe AI, ``Introducing System One Models \& Jev.''
  \url{https://typesafe.ai/blog/introducing-system-one-models-and-jev}
\end{thebibliography}

\section{Open problems}
(1) A privately-known $\pi$ for Jev requires the real API schema. (2) Measuring $H(\tau)$ needs
real agent traffic. (3) How to communicate with the agent vendor (Meta/Muse) is deferred.

\end{document}
